\documentclass[11pt]{article}
% Preámbulo común de los informes E1, E2 y E3, con aspecto de Word: Calibri,
% paleta del tema de Office, bandas y portada.
% Cada informe lo carga con % Preámbulo común de los informes E1, E2 y E3, con aspecto de Word: Calibri,
% paleta del tema de Office, bandas y portada.
% Cada informe lo carga con % Preámbulo común de los informes E1, E2 y E3, con aspecto de Word: Calibri,
% paleta del tema de Office, bandas y portada.
% Cada informe lo carga con \input{../../latex/preambulo.tex}.
%   \portada[Alumnos]{sigla y curso}{tipo de entrega}{título}{autores}{fecha}
\usepackage[spanish,es-tabla]{babel}
\usepackage[utf8]{inputenc}
\usepackage[T1]{fontenc}
% Carlito: mismas medidas que Calibri, la fuente de Word
\usepackage[sfdefault,lining]{carlito}
\usepackage[letterpaper,margin=2.2cm,top=2.8cm,bottom=2.8cm]{geometry}
\usepackage{xcolor}
\usepackage{graphicx}
\usepackage{array}
\usepackage{colortbl}
\usepackage{tabularx}
\usepackage{microtype}  % reparte el justificado en letras y espacios: menos huecos en columnas angostas
\usepackage{ragged2e}
\setlength{\JustifyingParindent}{0pt}  % sin sangría, como el resto del texto
\usepackage{enumitem}
\usepackage{tikz}
\usetikzlibrary{positioning,arrows.meta,calc,shapes}
\usepackage{eso-pic}
\usepackage{titlesec}
\usepackage{listings}
\usepackage[hidelinks]{hyperref}
\graphicspath{{../../assets/}}

% paleta del tema de Office
\definecolor{banda}{HTML}{4F81BD}
\definecolor{titulo}{HTML}{2E75B5}
\definecolor{subtitulo}{HTML}{5B9BD5}
\definecolor{pizarra}{HTML}{273540}
\definecolor{tablaencabezado}{HTML}{DEEAF6}  % el azul claro de los diagramas
\definecolor{fondosuave}{HTML}{DEEAF6}
\definecolor{grisoffice}{HTML}{7F7F7F}
\definecolor{vinotinto}{HTML}{7B1E2B}

% texto agregado o cambiado por el asistente, pendiente de revisión del equipo
\newcommand{\nuevo}[1]{\textcolor{vinotinto}{#1}}

% títulos como los estilos de título de Word
\titleformat{\section}{\large\bfseries\color{subtitulo}}{}{0pt}{}
\titleformat{\subsection}{\normalsize\bfseries\color{subtitulo}}{}{0pt}{}
\titlespacing*{\section}{0pt}{2ex}{1ex}
\titlespacing*{\subsection}{0pt}{1.5ex}{0.6ex}
\setcounter{secnumdepth}{0}

\setlist{itemsep=0.4ex,leftmargin=1.6em}
\setlist[itemize]{label=\textbullet}
\setlength{\parindent}{0pt}
\setlength{\parskip}{8pt plus 2pt}  % espacio después de párrafo de Word

% tablas tipo Word: grilla completa y encabezado con fondo
% celdas justificadas como el resto del texto; Y es la columna flexible de tabularx
\newcolumntype{L}[1]{>{\justifying\arraybackslash}p{#1}}
\newcolumntype{Y}{>{\justifying\arraybackslash}X}
\newcommand{\enc}[1]{\cellcolor{tablaencabezado}#1}
\newcommand{\iniciartabla}{\renewcommand{\arraystretch}{1.3}\arrayrulecolor{black!70}}

% notas al pie con números también dentro de bloques (minipage), como en Word
\renewcommand{\thempfootnote}{\arabic{mpfootnote}}

% pregunta que abre una respuesta
\newcommand{\pregunta}[1]{\textbf{\color{subtitulo}#1}}

% SQL y pseudocódigo
\lstset{basicstyle=\ttfamily\small,columns=fullflexible,keepspaces=true,
  frame=single,rulecolor=\color{grisoffice},backgroundcolor=\color{fondosuave},
  keywordstyle=\color{titulo}\bfseries,language=SQL,
  morekeywords={BEGIN,COMMIT,ROLLBACK},literate={á}{{\'a}}1 {é}{{\'e}}1
  {í}{{\'i}}1 {ó}{{\'o}}1 {ú}{{\'u}}1 {ñ}{{\~n}}1}

% portada: bandas completas; páginas interiores: media banda arriba a la
% izquierda, media banda abajo a la derecha con el número y logo de marca de agua
\pagestyle{empty}
\newif\ifenportada
\AddToShipoutPictureBG{%
  \ifenportada
    \AtPageUpperLeft{\color{banda}\raisebox{-1.2cm}{\rule{\paperwidth}{0.6cm}}}%
    \AtPageLowerLeft{\color{banda}\raisebox{0.6cm}{\rule{\paperwidth}{0.6cm}}}%
  \else
    \AtPageUpperLeft{\color{banda}\raisebox{-1.2cm}{\rule{0.62\paperwidth}{0.6cm}}}%
    \AtPageLowerLeft{%
      \put(\LenToUnit{0.45\paperwidth},\LenToUnit{0.6cm}){\color{banda}\rule{0.55\paperwidth}{0.6cm}}%
      \put(\LenToUnit{0.45\paperwidth},\LenToUnit{0.78cm}){\hspace{0.2cm}\color{white}\small\thepage}%
      \put(\LenToUnit{0.7cm},\LenToUnit{0.5cm}){\logomarcaagua}}%
  \fi}
\newcommand{\logomarcaagua}{\tikz\node[opacity=0.08,inner sep=0]{\includegraphics[width=2cm]{logo_uc.png}};}

% \portada[rótulo]{sigla y curso}{tipo de entrega}{título}{autores}{fecha}
\newcommand{\portada}[6][Alumno]{%
  \global\enportadatrue
  \begin{titlepage}
    \centering
    \vspace*{2.2cm}
    \includegraphics[width=8cm]{logo_uc.png}\par\vspace{1.6cm}
    {\large\bfseries\color{titulo} Magíster en Inteligencia Artificial - MIA\par}\vspace{0.4cm}
    {\large\bfseries\color{subtitulo} Instituto de Ingeniería Matemática y Computacional - IMC\par}\vspace{0.5cm}
    {\LARGE\bfseries\color{pizarra} #2\par}\vspace{0.5cm}
    {\itshape\bfseries\color{subtitulo} #3\par}\vspace{0.2cm}
    {\large #4\par}\vspace{0.2cm}
    {\color{banda}\rule{0.8\textwidth}{0.6pt}\par}
    \vspace{2cm}
    \raggedright
    {\textbf{#1:} #5\par}\vspace{0.3cm}
    {\textbf{Fecha de entrega:} #6\par}
    \vfill
  \end{titlepage}
  \global\enportadafalse}
.
%   \portada[Alumnos]{sigla y curso}{tipo de entrega}{título}{autores}{fecha}
\usepackage[spanish,es-tabla]{babel}
\usepackage[utf8]{inputenc}
\usepackage[T1]{fontenc}
% Carlito: mismas medidas que Calibri, la fuente de Word
\usepackage[sfdefault,lining]{carlito}
\usepackage[letterpaper,margin=2.2cm,top=2.8cm,bottom=2.8cm]{geometry}
\usepackage{xcolor}
\usepackage{graphicx}
\usepackage{array}
\usepackage{colortbl}
\usepackage{tabularx}
\usepackage{microtype}  % reparte el justificado en letras y espacios: menos huecos en columnas angostas
\usepackage{ragged2e}
\setlength{\JustifyingParindent}{0pt}  % sin sangría, como el resto del texto
\usepackage{enumitem}
\usepackage{tikz}
\usetikzlibrary{positioning,arrows.meta,calc,shapes}
\usepackage{eso-pic}
\usepackage{titlesec}
\usepackage{listings}
\usepackage[hidelinks]{hyperref}
\graphicspath{{../../assets/}}

% paleta del tema de Office
\definecolor{banda}{HTML}{4F81BD}
\definecolor{titulo}{HTML}{2E75B5}
\definecolor{subtitulo}{HTML}{5B9BD5}
\definecolor{pizarra}{HTML}{273540}
\definecolor{tablaencabezado}{HTML}{DEEAF6}  % el azul claro de los diagramas
\definecolor{fondosuave}{HTML}{DEEAF6}
\definecolor{grisoffice}{HTML}{7F7F7F}
\definecolor{vinotinto}{HTML}{7B1E2B}

% texto agregado o cambiado por el asistente, pendiente de revisión del equipo
\newcommand{\nuevo}[1]{\textcolor{vinotinto}{#1}}

% títulos como los estilos de título de Word
\titleformat{\section}{\large\bfseries\color{subtitulo}}{}{0pt}{}
\titleformat{\subsection}{\normalsize\bfseries\color{subtitulo}}{}{0pt}{}
\titlespacing*{\section}{0pt}{2ex}{1ex}
\titlespacing*{\subsection}{0pt}{1.5ex}{0.6ex}
\setcounter{secnumdepth}{0}

\setlist{itemsep=0.4ex,leftmargin=1.6em}
\setlist[itemize]{label=\textbullet}
\setlength{\parindent}{0pt}
\setlength{\parskip}{8pt plus 2pt}  % espacio después de párrafo de Word

% tablas tipo Word: grilla completa y encabezado con fondo
% celdas justificadas como el resto del texto; Y es la columna flexible de tabularx
\newcolumntype{L}[1]{>{\justifying\arraybackslash}p{#1}}
\newcolumntype{Y}{>{\justifying\arraybackslash}X}
\newcommand{\enc}[1]{\cellcolor{tablaencabezado}#1}
\newcommand{\iniciartabla}{\renewcommand{\arraystretch}{1.3}\arrayrulecolor{black!70}}

% notas al pie con números también dentro de bloques (minipage), como en Word
\renewcommand{\thempfootnote}{\arabic{mpfootnote}}

% pregunta que abre una respuesta
\newcommand{\pregunta}[1]{\textbf{\color{subtitulo}#1}}

% SQL y pseudocódigo
\lstset{basicstyle=\ttfamily\small,columns=fullflexible,keepspaces=true,
  frame=single,rulecolor=\color{grisoffice},backgroundcolor=\color{fondosuave},
  keywordstyle=\color{titulo}\bfseries,language=SQL,
  morekeywords={BEGIN,COMMIT,ROLLBACK},literate={á}{{\'a}}1 {é}{{\'e}}1
  {í}{{\'i}}1 {ó}{{\'o}}1 {ú}{{\'u}}1 {ñ}{{\~n}}1}

% portada: bandas completas; páginas interiores: media banda arriba a la
% izquierda, media banda abajo a la derecha con el número y logo de marca de agua
\pagestyle{empty}
\newif\ifenportada
\AddToShipoutPictureBG{%
  \ifenportada
    \AtPageUpperLeft{\color{banda}\raisebox{-1.2cm}{\rule{\paperwidth}{0.6cm}}}%
    \AtPageLowerLeft{\color{banda}\raisebox{0.6cm}{\rule{\paperwidth}{0.6cm}}}%
  \else
    \AtPageUpperLeft{\color{banda}\raisebox{-1.2cm}{\rule{0.62\paperwidth}{0.6cm}}}%
    \AtPageLowerLeft{%
      \put(\LenToUnit{0.45\paperwidth},\LenToUnit{0.6cm}){\color{banda}\rule{0.55\paperwidth}{0.6cm}}%
      \put(\LenToUnit{0.45\paperwidth},\LenToUnit{0.78cm}){\hspace{0.2cm}\color{white}\small\thepage}%
      \put(\LenToUnit{0.7cm},\LenToUnit{0.5cm}){\logomarcaagua}}%
  \fi}
\newcommand{\logomarcaagua}{\tikz\node[opacity=0.08,inner sep=0]{\includegraphics[width=2cm]{logo_uc.png}};}

% \portada[rótulo]{sigla y curso}{tipo de entrega}{título}{autores}{fecha}
\newcommand{\portada}[6][Alumno]{%
  \global\enportadatrue
  \begin{titlepage}
    \centering
    \vspace*{2.2cm}
    \includegraphics[width=8cm]{logo_uc.png}\par\vspace{1.6cm}
    {\large\bfseries\color{titulo} Magíster en Inteligencia Artificial - MIA\par}\vspace{0.4cm}
    {\large\bfseries\color{subtitulo} Instituto de Ingeniería Matemática y Computacional - IMC\par}\vspace{0.5cm}
    {\LARGE\bfseries\color{pizarra} #2\par}\vspace{0.5cm}
    {\itshape\bfseries\color{subtitulo} #3\par}\vspace{0.2cm}
    {\large #4\par}\vspace{0.2cm}
    {\color{banda}\rule{0.8\textwidth}{0.6pt}\par}
    \vspace{2cm}
    \raggedright
    {\textbf{#1:} #5\par}\vspace{0.3cm}
    {\textbf{Fecha de entrega:} #6\par}
    \vfill
  \end{titlepage}
  \global\enportadafalse}
.
%   \portada[Alumnos]{sigla y curso}{tipo de entrega}{título}{autores}{fecha}
\usepackage[spanish,es-tabla]{babel}
\usepackage[utf8]{inputenc}
\usepackage[T1]{fontenc}
% Carlito: mismas medidas que Calibri, la fuente de Word
\usepackage[sfdefault,lining]{carlito}
\usepackage[letterpaper,margin=2.2cm,top=2.8cm,bottom=2.8cm]{geometry}
\usepackage{xcolor}
\usepackage{graphicx}
\usepackage{array}
\usepackage{colortbl}
\usepackage{tabularx}
\usepackage{microtype}  % reparte el justificado en letras y espacios: menos huecos en columnas angostas
\usepackage{ragged2e}
\setlength{\JustifyingParindent}{0pt}  % sin sangría, como el resto del texto
\usepackage{enumitem}
\usepackage{tikz}
\usetikzlibrary{positioning,arrows.meta,calc,shapes}
\usepackage{eso-pic}
\usepackage{titlesec}
\usepackage{listings}
\usepackage[hidelinks]{hyperref}
\graphicspath{{../../assets/}}

% paleta del tema de Office
\definecolor{banda}{HTML}{4F81BD}
\definecolor{titulo}{HTML}{2E75B5}
\definecolor{subtitulo}{HTML}{5B9BD5}
\definecolor{pizarra}{HTML}{273540}
\definecolor{tablaencabezado}{HTML}{DEEAF6}  % el azul claro de los diagramas
\definecolor{fondosuave}{HTML}{DEEAF6}
\definecolor{grisoffice}{HTML}{7F7F7F}
\definecolor{vinotinto}{HTML}{7B1E2B}

% texto agregado o cambiado por el asistente, pendiente de revisión del equipo
\newcommand{\nuevo}[1]{\textcolor{vinotinto}{#1}}

% títulos como los estilos de título de Word
\titleformat{\section}{\large\bfseries\color{subtitulo}}{}{0pt}{}
\titleformat{\subsection}{\normalsize\bfseries\color{subtitulo}}{}{0pt}{}
\titlespacing*{\section}{0pt}{2ex}{1ex}
\titlespacing*{\subsection}{0pt}{1.5ex}{0.6ex}
\setcounter{secnumdepth}{0}

\setlist{itemsep=0.4ex,leftmargin=1.6em}
\setlist[itemize]{label=\textbullet}
\setlength{\parindent}{0pt}
\setlength{\parskip}{8pt plus 2pt}  % espacio después de párrafo de Word

% tablas tipo Word: grilla completa y encabezado con fondo
% celdas justificadas como el resto del texto; Y es la columna flexible de tabularx
\newcolumntype{L}[1]{>{\justifying\arraybackslash}p{#1}}
\newcolumntype{Y}{>{\justifying\arraybackslash}X}
\newcommand{\enc}[1]{\cellcolor{tablaencabezado}#1}
\newcommand{\iniciartabla}{\renewcommand{\arraystretch}{1.3}\arrayrulecolor{black!70}}

% notas al pie con números también dentro de bloques (minipage), como en Word
\renewcommand{\thempfootnote}{\arabic{mpfootnote}}

% pregunta que abre una respuesta
\newcommand{\pregunta}[1]{\textbf{\color{subtitulo}#1}}

% SQL y pseudocódigo
\lstset{basicstyle=\ttfamily\small,columns=fullflexible,keepspaces=true,
  frame=single,rulecolor=\color{grisoffice},backgroundcolor=\color{fondosuave},
  keywordstyle=\color{titulo}\bfseries,language=SQL,
  morekeywords={BEGIN,COMMIT,ROLLBACK},literate={á}{{\'a}}1 {é}{{\'e}}1
  {í}{{\'i}}1 {ó}{{\'o}}1 {ú}{{\'u}}1 {ñ}{{\~n}}1}

% portada: bandas completas; páginas interiores: media banda arriba a la
% izquierda, media banda abajo a la derecha con el número y logo de marca de agua
\pagestyle{empty}
\newif\ifenportada
\AddToShipoutPictureBG{%
  \ifenportada
    \AtPageUpperLeft{\color{banda}\raisebox{-1.2cm}{\rule{\paperwidth}{0.6cm}}}%
    \AtPageLowerLeft{\color{banda}\raisebox{0.6cm}{\rule{\paperwidth}{0.6cm}}}%
  \else
    \AtPageUpperLeft{\color{banda}\raisebox{-1.2cm}{\rule{0.62\paperwidth}{0.6cm}}}%
    \AtPageLowerLeft{%
      \put(\LenToUnit{0.45\paperwidth},\LenToUnit{0.6cm}){\color{banda}\rule{0.55\paperwidth}{0.6cm}}%
      \put(\LenToUnit{0.45\paperwidth},\LenToUnit{0.78cm}){\hspace{0.2cm}\color{white}\small\thepage}%
      \put(\LenToUnit{0.7cm},\LenToUnit{0.5cm}){\logomarcaagua}}%
  \fi}
\newcommand{\logomarcaagua}{\tikz\node[opacity=0.08,inner sep=0]{\includegraphics[width=2cm]{logo_uc.png}};}

% \portada[rótulo]{sigla y curso}{tipo de entrega}{título}{autores}{fecha}
\newcommand{\portada}[6][Alumno]{%
  \global\enportadatrue
  \begin{titlepage}
    \centering
    \vspace*{2.2cm}
    \includegraphics[width=8cm]{logo_uc.png}\par\vspace{1.6cm}
    {\large\bfseries\color{titulo} Magíster en Inteligencia Artificial - MIA\par}\vspace{0.4cm}
    {\large\bfseries\color{subtitulo} Instituto de Ingeniería Matemática y Computacional - IMC\par}\vspace{0.5cm}
    {\LARGE\bfseries\color{pizarra} #2\par}\vspace{0.5cm}
    {\itshape\bfseries\color{subtitulo} #3\par}\vspace{0.2cm}
    {\large #4\par}\vspace{0.2cm}
    {\color{banda}\rule{0.8\textwidth}{0.6pt}\par}
    \vspace{2cm}
    \raggedright
    {\textbf{#1:} #5\par}\vspace{0.3cm}
    {\textbf{Fecha de entrega:} #6\par}
    \vfill
  \end{titlepage}
  \global\enportadafalse}


\newcolumntype{Y}{>{\raggedright\arraybackslash}X}

\begin{document}

\portada[Alumnos]{INF3862-1 Plataformas Big Data para Data Science}
  {Proyecto del curso, Entrega 1: Diseño de datos}
  {Hábitos de actividad y alimentación asociados a mejores indicadores de salud}
  {Miguel Angel Cervantes Buleje y Roger Ruben Reyes Miranda}
  {16 de noviembre de 2026}

\section{4.1 Objetivo final del proyecto}

El proyecto compara los hábitos de actividad y de alimentación de los usuarios
de la aplicación según su frecuencia cardíaca en reposo, el indicador del reloj
que se usa como referencia de condición física. Los resultados describen
asociaciones en los datos, no causas ni diagnósticos.

\begingroup\iniciartabla\small
\begin{tabularx}{\textwidth}{|L{2.3cm}|Y|}\hline
\enc{\textbf{Pregunta}} & \enc{\textbf{Respuesta}}\\\hline
Resultado & Un dashboard para el equipo de nutrición y contenidos de la
aplicación. Agrupa a los usuarios en cuartiles de frecuencia cardíaca en reposo
y muestra, por cuartil, pasos diarios, minutos de actividad y nutrientes de las
recetas que preparan. Apoya la decisión de qué recetas y actividades promover
en las recomendaciones.\\\hline
Datos & Del reloj: 315 relojes de 300 usuarios, 15540132 mediciones (115.5 MB
comprimidos), 4840 sesiones y 22106 comidas registradas en 61 días. De
Food.com: 522517 recetas con su información nutricional (178.7 MB en Parquet).
Las mediciones llegan en lotes cuando el reloj se sincroniza; las comidas y
las vinculaciones de relojes, como transacciones de la aplicación.\\\hline
Operaciones & Unificar los formatos de los dos modelos de reloj, quitar
mediciones repetidas, calcular por usuario y día los pasos, los minutos de
sesión y la frecuencia cardíaca en reposo, unir cada comida con la nutrición de
su receta y agregar por cuartil. No se predice: el objetivo es descriptivo.\\\hline
Exigencias & El dashboard se actualiza una vez al día; un resultado puede
esperar horas. Se toleran lotes que llegan tarde si se recalculan los días que
afectan. No se toleran mediciones duplicadas, porque suman pasos dos veces. Una
frecuencia cardíaca nula se excluye del cálculo, no se rellena.\\\hline
Recursos & Un equipo de datos pequeño y presupuesto bajo. Los datos del
proyecto suman cerca de 296 MB y caben en una máquina; el almacenamiento de
objetos guarda los archivos originales y las tablas analíticas.\\\hline
\end{tabularx}
\endgroup

\textbf{Frecuencia cardíaca en reposo:} promedio de la frecuencia cardíaca entre
las 2:00 y las 5:00, en la hora local del usuario y fuera de sesiones de
actividad. Mide el pulso sin esfuerzo; un valor más bajo se interpreta como mejor
condición física, sin valor diagnóstico.

\section{4.2 Fuentes y recepción de datos}

\begingroup\iniciartabla\footnotesize
\begin{tabularx}{\textwidth}{|L{1.6cm}|L{2.6cm}|Y|Y|Y|}\hline
\enc{\textbf{Fuente}} & \enc{\textbf{Registro y tipo}} &
\enc{\textbf{Volumen y crecimiento}} & \enc{\textbf{Recepción}} &
\enc{\textbf{Problemas esperados}}\\\hline
Recetas & Una receta. Estructurado, con columnas de listas &
522517 filas, 178.7 MB. Unas 70 mil nuevas al año (70204 publicadas en 2007) &
Transacción en línea al publicar; al análisis, batch diario &
Duraciones como texto (PT45M), recetas editadas\\\hline
Comidas & Una receta preparada por un usuario. Estructurado &
22106 filas, 1.0 MB. Unas 362 por día con 300 usuarios &
Transacción en línea en segundos; al análisis, batch diario &
Recetas fuera del catálogo (48 comidas), doble envío\\\hline
Dispositivos & Un reloj vinculado a un usuario. Estructurado &
315 relojes de 300 usuarios, 15 desvinculados &
Transacción en línea al vincular o desvincular &
Usuarios con más de un reloj\\\hline
Sesiones & Una actividad. Semiestructurado (JSON) &
4840 filas, 0.8 MB. Unas 79 por día &
Lotes en cada sincronización; latencia de horas &
Distancia nula en fuerza, repetidas por reintentos\\\hline
Mediciones & Una medición del reloj. Semiestructurado (JSON comprimido) &
15540132 filas, 115.5 MB. Unas 254756 por día (1.9 MB) &
Lotes por sincronización, por hora de recepción; latencia de horas &
Lotes atrasados, repetidas, frecuencia cardíaca nula (2.3\%), dos formatos,
muestreo de 1 min, 5 min o 10 s\\\hline
\end{tabularx}
\endgroup

\textbf{¿Se interpreta el esquema al escribir o al leer, y por qué?} Al escribir
en las fuentes que crea la aplicación (recetas, comidas, dispositivos): las
reglas de integridad se validan en la transacción y un registro inválido se
rechaza antes de afectar la operación. Al leer en las que llegan del reloj
(sesiones y mediciones): Pulso S envía \texttt{frecuencia\_cardiaca} y fecha ISO
8601, y FitBand 2 envía \texttt{fc} y milisegundos desde 1970. Guardar el lote
original y mapear los nombres al procesarlo permite reprocesar si cambia una
regla o aparece un modelo nuevo, sin perder lo recibido.

\section{4.3 Estructura de los datos}

\subsection{a) Nivel operacional (OLTP)}

\begin{itemize}
  \item \textbf{Usuario} (\texttt{usuario\_id} PK, nombre, zona\_horaria),
  \textbf{receta} (\texttt{receta\_id} PK, \texttt{autor\_id} FK a usuario,
  nombre, categoría, porciones, calorías, grasa saturada, sodio, azúcar, fibra,
  proteína) y \textbf{comida} (\texttt{comida\_id} PK, \texttt{usuario\_id} FK,
  \texttt{receta\_id} FK, fecha\_hora, porciones). Modelo relacional: la
  aplicación lee y escribe pocas filas por clave, y las claves foráneas impiden
  registrar una comida de una receta inexistente.
  \item \textbf{Reloj} (\texttt{device\_id} PK, modelo), \textbf{vinculación}
  (\texttt{vinculacion\_id} PK, \texttt{device\_id} FK, \texttt{usuario\_id} FK,
  desde, hasta) y \textbf{sesión} (\texttt{session\_id} PK, \texttt{device\_id}
  FK, tipo, inicio, fin, distancia\_km, calorías). Modelo relacional: la
  vinculación necesita una transacción y la aplicación lista las sesiones de un
  usuario.
  \item \textbf{Medición} (\texttt{device\_id} y \texttt{ts\_medicion} como PK
  compuesta, ts\_recepcion, frecuencia\_cardiaca, pasos, calorías). Modelo de
  series temporales: llegan muchos eventos por reloj y se consultan por reloj y
  rango de tiempo. La PK compuesta descarta la medición repetida de un
  reintento.
\end{itemize}

\textbf{¿Qué regla de integridad exige una transacción?} Un reloj está
vinculado a un solo usuario a la vez: al vincularlo a otro usuario, cerrar el
vínculo anterior y abrir el nuevo son una unidad (consulta 1 del anexo). La figura~\ref{fig:er} muestra el diagrama entidad-relación.

\subsection{b) Nivel analítico (OLAP)}

\begin{itemize}
  \item \textbf{Hecho \texttt{hecho\_usuario\_dia}:} una fila por usuario y día
  local. Medidas: pasos, minutos de sesión, calorías gastadas, frecuencia
  cardíaca en reposo, número de comidas, y calorías, sodio, grasa saturada,
  azúcar, fibra y proteína consumidas. Claves: \texttt{usuario\_key},
  \texttt{fecha\_key}, \texttt{reloj\_key}.
  \item \textbf{Hecho \texttt{hecho\_comida}:} una fila por comida registrada,
  con porciones y nutrientes consumidos; claves a usuario, receta y fecha.
  \item \textbf{Dimensiones:} fecha (día, día de la semana, semana, mes),
  usuario (zona horaria), reloj (modelo) y receta (nombre, categoría,
  nutrientes por porción).
  \item \textbf{Supuesto:} la información nutricional de Food.com es por
  porción; los nutrientes consumidos son los de la receta multiplicados por las
  porciones de la comida.
\end{itemize}

\textbf{¿Dónde viven los datos crudos y los analíticos, y en qué formato?} En un
Lakehouse. Los lotes originales quedan en almacenamiento de objetos tal como
llegan (JSON comprimido y CSV), para reprocesarlos con schema-on-read. Las tablas
analíticas son tablas Delta (archivos Parquet y un registro de transacciones)
sobre el mismo almacenamiento: cuando un lote atrasado obliga a reescribir un
día, el cambio se confirma como una unidad y el dashboard nunca lee un día a
medio reescribir. Un Warehouse guarda datos ya integrados y dejaría fuera los
JSON originales que se necesitan para reprocesar.

\textbf{¿Qué organización física y partición se proponen?} Por columnas
(Parquet): el dashboard lee pocas medidas de muchas filas, y el formato
columnar lee solo esas columnas y las comprime. Las mediciones limpias se
particionan por fecha de medición, de modo que recalcular un día afectado por
un lote atrasado lee solo ese día. Los hechos se particionan por mes, y las
consultas con rango de fechas podan los meses que no piden.

\section{4.4 Consultas de ejemplo}

El anexo A detalla cada consulta con su pseudocódigo, justificación y lo que la
haría rápida.

\begingroup\iniciartabla\small
\begin{tabularx}{\textwidth}{|c|Y|L{2cm}|L{3.3cm}|}\hline
\enc{\textbf{N.}} & \enc{\textbf{Propósito y usuario}} &
\enc{\textbf{Clasificación}} & \enc{\textbf{Qué la haría rápida}}\\\hline
1 & Vincular un reloj a un usuario (la aplicación) & OLTP, escribe &
Índice por \texttt{device\_id}\\\hline
2 & Registrar una comida (la aplicación) & OLTP, escribe & Índices de las
claves foráneas\\\hline
3 & Listar las comidas de hoy de un usuario (la aplicación) & OLTP, lee &
Índice por usuario y fecha\\\hline
4 & Comparar hábitos por cuartil de frecuencia cardíaca en reposo (el equipo
de nutrición) & OLAP & Columnar y partición por mes\\\hline
5 & Pasos semanales por modelo de reloj (un analista) & OLAP & Columnar y
partición por mes\\\hline
6 & Recetas más preparadas por cuartil (el equipo de nutrición) & OLAP &
Columnar\\\hline
7 & Resumen de hoy de un usuario en la aplicación & Frontera & Resumen
precalculado\\\hline
\end{tabularx}
\endgroup

\textbf{¿Qué pasaría sin transacción en la consulta 1?} Si la escritura falla
después de cerrar el vínculo anterior, el reloj queda sin usuario y sus
mediciones no se asignan a nadie. Si dos vinculaciones del mismo reloj corren a
la vez, ambas pueden ver el reloj libre y dejar dos vínculos abiertos.

\textbf{¿Cómo se clasifica la consulta 7 y dónde se ejecuta?} Agrega cientos de
mediciones del día, como una consulta OLAP, pero la ejecuta la aplicación para
un solo usuario, con todos los usuarios activos a la vez y latencia de segundos, como
una OLTP. Se clasifica como OLTP sobre un resumen precalculado: la base de
series temporales actualiza un acumulado por reloj y día al recibir cada lote,
y la aplicación lee una fila. Así no recorre mediciones en la base relacional
ni compite con las comidas y vinculaciones.

\section{4.5 Arquitectura inicial}

La figura~\ref{fig:arquitectura} muestra el recorrido de los datos. La carga al almacenamiento
analítico es batch y diaria, porque el dashboard se actualiza una vez al día.

\begin{figure}[tb]
\centering
\begin{tikzpicture}[
  caja/.style={draw=banda,fill=fondosuave,rounded corners=2pt,align=center,
    font=\footnotesize,minimum height=1.1cm,text width=2.6cm},
  titulo/.style={font=\footnotesize\bfseries\color{titulo}},
  flecha/.style={-{Stealth[length=2mm]},draw=pizarra}]
\node[caja] (app) {Aplicación\\recetas, comidas, vínculos};
\node[caja,below=1.2cm of app] (reloj) {Relojes\\mediciones, sesiones};
\node[caja,right=0.6cm of app] (api) {Transacciones en línea};
\node[caja,right=0.6cm of reloj] (lotes) {Recepción de lotes de sincronización};
\node[caja,right=0.6cm of api] (rel) {Base relacional\\usuarios, recetas, comidas, vínculos, sesiones};
\node[caja,right=0.6cm of lotes] (st) {Base de series temporales\\mediciones y acumulado diario};
\node[caja,right=0.6cm of rel] (delta) {Tablas Delta\\hechos y dimensiones};
\node[caja,right=0.6cm of st] (obj) {Almacenamiento de objetos\\lotes originales y extractos};
\node[caja,right=0.6cm of delta,yshift=0.4cm] (dash) {Dashboard del equipo de nutrición};
\node[caja,below=0.3cm of dash] (nb) {Notebooks de análisis};
\node[caja,right=0.6cm of obj] (appc) {Aplicación\\resumen del día};
\node[titulo,above=0.3cm of app] {Fuentes};
\node[titulo,above=0.3cm of api] {Recepción};
\node[titulo,above=0.3cm of rel] {Operacional};
\node[titulo,above=0.3cm of delta] {Analítico};
\node[titulo,above=0.3cm of dash] {Consumo};
\draw[flecha] (app) -- (api);
\draw[flecha] (reloj) -- (lotes);
\draw[flecha] (api) -- (rel);
\draw[flecha] (lotes) -- (st);
\draw[flecha] (lotes.south) |- ([yshift=-0.4cm]obj.south) -- (obj.south);
\draw[flecha] (rel.south east) -- node[right,font=\scriptsize,pos=0.3]{batch diario} (obj.north west);
\draw[flecha] (obj) -- (delta);
\draw[flecha] (delta) -- (dash);
\draw[flecha] (delta) -- (nb);
\draw[flecha] (st.south) |- ([yshift=-0.8cm]appc.south) -- (appc.south);
\end{tikzpicture}
\caption{Arquitectura inicial por capacidades.}\label{fig:arquitectura}
\end{figure}

\begin{figure}[tb]
\centering
\begin{tikzpicture}[
  ent/.style={draw=banda,fill=fondosuave,align=left,font=\scriptsize,
    inner sep=3pt},
  rel/.style={draw=pizarra}]
\node[ent] (usu) {\textbf{usuario}\\usuario\_id PK\\nombre\\zona\_horaria};
\node[ent,right=1.6cm of usu] (com) {\textbf{comida}\\comida\_id PK\\usuario\_id FK\\receta\_id FK\\fecha\_hora, porciones};
\node[ent,right=1.6cm of com] (rec) {\textbf{receta}\\receta\_id PK\\autor\_id FK\\nombre, categoría\\nutrientes por porción};
\node[ent,below=0.9cm of usu] (vin) {\textbf{vinculación}\\vinculacion\_id PK\\device\_id FK\\usuario\_id FK\\desde, hasta};
\node[ent,right=1.6cm of vin] (rel) {\textbf{reloj}\\device\_id PK\\modelo};
\node[ent,right=1.6cm of rel] (ses) {\textbf{sesión}\\session\_id PK\\device\_id FK\\tipo, inicio, fin};
\node[ent,below=0.6cm of rel] (med) {\textbf{medición} (series temporales)\\device\_id, ts\_medicion PK\\frecuencia\_cardiaca, pasos};
\draw[rel] (usu) -- node[above,font=\scriptsize]{1 : N} (com);
\draw[rel] (com) -- node[above,font=\scriptsize]{N : 1} (rec);
\draw[rel] (usu) -- node[left,font=\scriptsize]{1 : N} (vin);
\draw[rel] (vin) -- node[above,font=\scriptsize]{N : 1} (rel);
\draw[rel] (rel) -- node[above,font=\scriptsize]{1 : N} (ses);
\draw[rel] (rel) -- node[right,font=\scriptsize]{1 : N} (med);
\draw[rel] (usu.north) -- ++(0,0.4) -| node[pos=0.25,above,font=\scriptsize]{autor 1 : N} (rec.north);
\end{tikzpicture}
\caption{Diagrama entidad-relación del nivel operacional (4.3 a).}\label{fig:er}
\end{figure}

\clearpage
\section{Anexo A. Consultas de ejemplo}

\par\noindent\begin{minipage}{\linewidth}\setlength{\parskip}{1ex}
\textbf{Consulta 1. Vincular un reloj a un usuario.} La ejecuta la aplicación
cuando el usuario vincula su reloj. OLTP: escribe dos filas y lee una por
clave, pocas veces al día, con latencia de milisegundos y muchos usuarios a la
vez. Rápida con un índice por \texttt{device\_id} en vinculación.
\begin{lstlisting}
BEGIN;
SELECT vinculacion_id FROM vinculacion
 WHERE device_id = :reloj AND hasta IS NULL
 FOR UPDATE;                         -- bloquea el vínculo abierto
UPDATE vinculacion SET hasta = now()
 WHERE device_id = :reloj AND hasta IS NULL;
INSERT INTO vinculacion (device_id, usuario_id, desde, hasta)
 VALUES (:reloj, :usuario, now(), NULL);
COMMIT;                              -- ante un error: ROLLBACK
\end{lstlisting}\end{minipage}\par\medskip

\par\noindent\begin{minipage}{\linewidth}\setlength{\parskip}{1ex}
\textbf{Consulta 2. Registrar una comida.} La ejecuta la aplicación cuando el
usuario marca una receta como preparada. OLTP: escribe una fila, unas 362 veces
al día con 300 usuarios, con latencia de milisegundos. Las claves foráneas
rechazan una receta inexistente. Rápida con índices en las claves foráneas.
\begin{lstlisting}
INSERT INTO comida (usuario_id, receta_id, fecha_hora, porciones)
 VALUES (:usuario, :receta, now(), :porciones);
\end{lstlisting}\end{minipage}\par\medskip

\par\noindent\begin{minipage}{\linewidth}\setlength{\parskip}{1ex}
\textbf{Consulta 3. Comidas de hoy de un usuario.} La ejecuta la aplicación al
abrir la pantalla del día. OLTP: lee unas pocas filas de un usuario, muchas
veces al día y con muchos usuarios a la vez, con latencia de milisegundos.
Rápida con un índice por (\texttt{usuario\_id}, \texttt{fecha\_hora}).
\begin{lstlisting}
SELECT r.nombre, c.porciones, r.calorias * c.porciones AS calorias
FROM comida c JOIN receta r ON r.receta_id = c.receta_id
WHERE c.usuario_id = :usuario AND c.fecha_hora >= :inicio_del_dia;
\end{lstlisting}\end{minipage}\par\medskip

\par\noindent\begin{minipage}{\linewidth}\setlength{\parskip}{1ex}
\textbf{Consulta 4. Hábitos por cuartil de frecuencia cardíaca en reposo.}
Responde el objetivo de 4.1; la ejecuta el dashboard del equipo de nutrición.
OLAP: lee todas las filas de \texttt{hecho\_usuario\_dia} del período (hasta 18300
con 300 usuarios y 61 días) y agrega, una vez al día, sin exigencia de
segundos. Rápida con almacenamiento columnar, que lee solo las medidas usadas,
y partición por mes, que poda los meses fuera del período.
\begin{lstlisting}
WITH perfil AS (
  SELECT usuario_key, AVG(fc_reposo) AS fc_reposo, AVG(pasos) AS pasos,
         AVG(minutos_sesion) AS minutos, AVG(sodio) AS sodio,
         AVG(grasa_saturada) AS grasa_sat, AVG(fibra) AS fibra
  FROM hecho_usuario_dia
  WHERE fecha_key BETWEEN :desde AND :hasta
  GROUP BY usuario_key)
SELECT NTILE(4) OVER (ORDER BY fc_reposo) AS cuartil, AVG(pasos),
       AVG(minutos), AVG(sodio), AVG(grasa_sat), AVG(fibra)
FROM perfil GROUP BY cuartil;
\end{lstlisting}\end{minipage}\par\medskip

\par\noindent\begin{minipage}{\linewidth}\setlength{\parskip}{1ex}
\textbf{Consulta 5. Pasos semanales por modelo de reloj.} La ejecuta un
analista para revisar si los dos modelos, que miden con distinta frecuencia,
reportan pasos comparables. OLAP: lee todo el período y agrupa por semana y
modelo, de forma ocasional. Rápida con almacenamiento columnar y partición por
mes.
\begin{lstlisting}
SELECT f.semana, r.modelo, AVG(h.pasos) AS pasos_promedio
FROM hecho_usuario_dia h
JOIN dim_fecha f ON f.fecha_key = h.fecha_key
JOIN dim_reloj r ON r.reloj_key = h.reloj_key
GROUP BY f.semana, r.modelo;
\end{lstlisting}\end{minipage}\par\medskip

\par\noindent\begin{minipage}{\linewidth}\setlength{\parskip}{1ex}
\textbf{Consulta 6. Recetas más preparadas por cuartil.} La ejecuta el equipo
de nutrición para elegir qué recetas promover. OLAP: lee todas las comidas del
período (22106 en 61 días) y agrupa, una vez al día. Rápida con almacenamiento
columnar.
\begin{lstlisting}
SELECT p.cuartil, d.nombre, COUNT(*) AS veces
FROM hecho_comida h
JOIN dim_receta d ON d.receta_key = h.receta_key
JOIN cuartil_usuario p ON p.usuario_key = h.usuario_key  -- de la consulta 4
GROUP BY p.cuartil, d.nombre
ORDER BY p.cuartil, veces DESC;      -- las 10 primeras de cada cuartil
\end{lstlisting}\end{minipage}\par\medskip

\par\noindent\begin{minipage}{\linewidth}\setlength{\parskip}{1ex}
\textbf{Consulta 7. Resumen de hoy de un usuario (frontera).} La ejecuta la
aplicación al abrir la pantalla del día. Agrega las mediciones del día, una por
minuto con Pulso S, pero para un solo usuario, muchas veces al día y
con latencia de segundos. Rápida con un resumen precalculado: el acumulado por
reloj y día se actualiza al recibir cada lote.
\begin{lstlisting}
-- al recibir un lote (base de series temporales)
UPSERT acumulado_dia (device_id, fecha, pasos, calorias)
  sumando pasos y calorias del lote sin repetidos;
-- en la aplicación
SELECT pasos, calorias FROM acumulado_dia
WHERE device_id = :reloj AND fecha = :hoy;
\end{lstlisting}\end{minipage}\par\medskip

\end{document}
